%! Logarithmic Functions — Unit 5, Lesson 5.3 — Slides (Beamer)
\documentclass[aspectratio=169, 11pt]{beamer}
\usepackage{precalculus-beamer}   % provides \navyheader{} and \sectionlabel[color]{}

\begin{document}

% TITLE SLIDE
{
\setbeamercolor{background canvas}{bg=navy}
\begin{frame}[plain]
  \vspace{2.8cm}\hspace{0.55cm}%
  \begin{minipage}{0.9\textwidth}
    {\color{goldacc}\bfseries\small LESSON 2.11}\\[0.5cm]
    {\color{white}\bfseries\LARGE Logarithmic Functions}\\[1.2cm]
    {\color{skymid}\small Precalculus~$\cdot$~Unit 2}
  \end{minipage}
\end{frame}
}

% SLIDE 1: Hook — Richter Scale
\begin{frame}[t, plain]
  \navyheader{Hook: The Richter Scale}
  \begin{block}{Richter Scale}
    Earthquake magnitude: $M = \log_{10}(I)$, where $I$ = intensity
  \end{block}

  \vspace{0.3cm}
  \begin{columns}[T]
    \begin{column}{0.5\textwidth}
      \begin{tabular}{|c|c|}
        \hline
        \textbf{Intensity $I$} & \textbf{Magnitude $M$} \\
        \hline
        10 & 1 \\
        100 & 2 \\
        1{,}000 & 3 \\
        10{,}000 & 4 \\
        \hline
      \end{tabular}
    \end{column}
    \begin{column}{0.46\textwidth}
      \begin{itemize}
        \setlength{\itemsep}{8pt}
        \item Each time $I$ multiplies by \textbf{10}\ldots
        \item \ldots $M$ increases by \textbf{1}
        \item What inputs $I$ make sense?
        \item What type of function is this?
      \end{itemize}
    \end{column}
  \end{columns}
\end{frame}

% SLIDE 2: Domain, Range, Asymptote
\begin{frame}[t, plain]
  \navyheader{Domain, Range, and Vertical Asymptote (EK 2.11.A.1, A.4)}

  \begin{block}{General Form: $f(x) = a\log_b x$, \; $b > 0$, $b \ne 1$, $a \ne 0$}
  \end{block}

  \vspace{0.25cm}
  \begin{columns}[T]
    \begin{column}{0.48\textwidth}
      \begin{itemize}
        \setlength{\itemsep}{10pt}
        \item \textbf{Domain:} $x > 0$, i.e., $(0, \infty)$
        \item \textbf{Range:} all real numbers $(-\infty, \infty)$
        \item \textbf{Vertical asymptote:} $x = 0$
      \end{itemize}
    \end{column}
    \begin{column}{0.48\textwidth}
      \begin{itemize}
        \setlength{\itemsep}{10pt}
        \item $\log_b(0)$ is \textbf{undefined}
        \item $\log_b(\text{negative})$ is \textbf{undefined}
        \item No $y$-intercept --- the graph never touches $x = 0$
      \end{itemize}
    \end{column}
  \end{columns}

  \vspace{0.3cm}
  \begin{block}{End Behavior (EK 2.11.A.4)}
    For $b > 1$, $a > 0$:\quad
    $\displaystyle\lim_{x \to 0^+} a\log_b x = -\infty$
    \qquad
    $\displaystyle\lim_{x \to \infty} a\log_b x = +\infty$
  \end{block}
\end{frame}

% SLIDE 3: Monotonicity and Concavity
\begin{frame}[t, plain]
  \navyheader{Monotonicity and Concavity (EK 2.11.A.2)}

  \begin{block}{From the Inverse Relationship}
    $f(x) = \log_b x$ is the inverse of $g(x) = b^x$, so it inherits structural properties.
  \end{block}

  \vspace{0.3cm}
  \begin{columns}[T]
    \begin{column}{0.48\textwidth}
      \textbf{When $b > 1$:}
      \begin{itemize}
        \setlength{\itemsep}{7pt}
        \item $g(x) = b^x$ is \textbf{always increasing}
        \item So $f(x) = \log_b x$ is \textbf{always increasing}
        \item Concave \textbf{down}
      \end{itemize}

      \vspace{0.2cm}
      \textbf{When $0 < b < 1$:}
      \begin{itemize}
        \setlength{\itemsep}{7pt}
        \item $g(x) = b^x$ is \textbf{always decreasing}
        \item So $f(x) = \log_b x$ is \textbf{always decreasing}
        \item Concave \textbf{up}
      \end{itemize}
    \end{column}
    \begin{column}{0.48\textwidth}
      \begin{block}{Key Consequence}
        \begin{itemize}
          \setlength{\itemsep}{8pt}
          \item \textbf{No local maximum or minimum} on the natural domain
          \item \textbf{No inflection point}
          \item The graph is either always concave up or always concave down
        \end{itemize}
      \end{block}
    \end{column}
  \end{columns}
\end{frame}

% SLIDE 4: End Behavior Cases
\begin{frame}[t, plain]
  \navyheader{End Behavior — All Cases (EK 2.11.A.4)}

  \vspace{0.2cm}
  \begin{center}
  \renewcommand{\arraystretch}{1.7}
  \begin{tabular}{|l|c|c|}
    \hline
    \textbf{Conditions} &
    $\displaystyle\lim_{x \to 0^+} a\log_b x$ &
    $\displaystyle\lim_{x \to \infty} a\log_b x$ \\
    \hline
    $b > 1$, $a > 0$ & $-\infty$ & $+\infty$ \\
    \hline
    $b > 1$, $a < 0$ & $+\infty$ & $-\infty$ \\
    \hline
    $0 < b < 1$, $a > 0$ & $+\infty$ & $-\infty$ \\
    \hline
    $0 < b < 1$, $a < 0$ & $-\infty$ & $+\infty$ \\
    \hline
  \end{tabular}
  \end{center}

  \vspace{0.2cm}
  \begin{block}{Memory Aid}
    Base $b > 1$ + positive $a$ \;\textrightarrow\; increasing log \;\textrightarrow\;
    $-\infty$ at left, $+\infty$ at right. \\
    Flip $a$ or flip $b$ \;\textrightarrow\; flip the end behavior direction.
  \end{block}
\end{frame}

% SLIDE 5: Additive Transformation Identification Test
\begin{frame}[t, plain]
  \navyheader{Additive Transformation Identification Test (EK 2.11.A.3)}

  \begin{block}{The Test}
    If the \textbf{inputs} of $f(x)$ are proportional over equal-length output intervals,
    then $f$ is logarithmic.
  \end{block}

  \vspace{0.25cm}
  \begin{columns}[T]
    \begin{column}{0.45\textwidth}
      \textbf{Example:} $f(x) = \log_3 x$

      \vspace{0.15cm}
      \begin{tabular}{|c|c|c|}
        \hline
        \textbf{$x$} & \textbf{$f(x)$} & \textbf{Ratio} \\
        \hline
        1  & 0 & --- \\
        3  & 1 & $\times 3$ \\
        9  & 2 & $\times 3$ \\
        27 & 3 & $\times 3$ \\
        81 & 4 & $\times 3$ \\
        \hline
      \end{tabular}

      \vspace{0.1cm}
      \small Inputs multiply by 3 each time output adds 1. \;\textbf{Logarithmic!}
    \end{column}
    \begin{column}{0.50\textwidth}
      \begin{block}{Key Nuance (EK 2.11.A.3)}
        \begin{itemize}
          \setlength{\itemsep}{7pt}
          \item $f$ logarithmic $\Rightarrow$ inputs proportional
          \item $g(x) = f(x+k)$ (additive shift) does \textbf{NOT} have proportional inputs
          \item If $g(x) = f(x+k)$ has proportional inputs, then $f$ is logarithmic
        \end{itemize}
      \end{block}
    \end{column}
  \end{columns}
\end{frame}

% SLIDE 6: Class Practice
\begin{frame}[t, plain]
  \navyheader{Class Practice}

  \begin{columns}[T]
    \begin{column}{0.48\textwidth}
      \textbf{1.} State all characteristics of $f(x) = 2\log_3 x$:
      \begin{itemize}
        \setlength{\itemsep}{6pt}
        \item Domain: \underline{\hspace{2.5cm}}
        \item Range: \underline{\hspace{2.5cm}}
        \item Asymptote: \underline{\hspace{2cm}}
        \item $\displaystyle\lim_{x\to 0^+} = $ \underline{\hspace{1.5cm}}
        \item $\displaystyle\lim_{x\to\infty} = $ \underline{\hspace{1.5cm}}
        \item Increasing or decreasing?
        \item Concave up or down?
      \end{itemize}
    \end{column}
    \begin{column}{0.48\textwidth}
      \textbf{2.} A table shows outputs $0, 1, 2, 3$ for a mystery function $h$ with
      inputs $4, 16, 64, 256$.

      \vspace{0.2cm}
      Is $h$ logarithmic? If so, what base?

      \vspace{0.8cm}
      \textbf{3.} True or false: $f(x) = \log_5 x$ has a $y$-intercept. Explain.
    \end{column}
  \end{columns}
\end{frame}

\end{document}
